\section*{Chapter \ref{ch:g9:ions} --- Ions and Ionic Solutions}

\begin{solution}{exo:g9:ions:1}
An \omterm{def:g7:atoms-and-molecules:atom}{atom} or group of \omterm{def:g7:atoms-and-molecules:atom}{atoms} that has lost or gained \omterm{def:g9:inside-the-atom:nucleus}{electrons} and carries a
charge. A \omterm{def:g9:ions:ion}{cation} is positive (\omterm{def:g9:inside-the-atom:nucleus}{electrons} lost), an \omterm{def:g9:ions:ion}{anion} negative
(\omterm{def:g9:inside-the-atom:nucleus}{electrons} gained).
\end{solution}

\begin{solution}{exo:g9:ions:2}
\ce{Mg^{2+}}: 12 \omterm{def:g9:inside-the-atom:proton}{protons}, 10 \omterm{def:g9:inside-the-atom:nucleus}{electrons}.
\end{solution}

\begin{solution}{exo:g9:ions:3}
\ce{KCl}, \ce{MgCl2}, \ce{CaO}.
\end{solution}

\begin{solution}{exo:g9:ions:4}
Salt water contains \omterm{def:g9:ions:ion}{ions}, \ce{Na+(aq)} and \ce{Cl-(aq)}, which move and
carry charge. Sugar \omterm{def:g7:atoms-and-molecules:molecule}{molecules} carry no charge.
\end{solution}

\begin{solution}{exo:g9:ions:5}
\ce{CuSO4}, \ce{AlCl3}, \ce{Na2CO3}.
\end{solution}

\begin{solution}{exo:g9:ions:6}
Iron(III), \ce{Fe^{3+}}: \ce{Fe^{3+}(aq) + 3OH-(aq) -> Fe(OH)3(s)}.
\end{solution}

\begin{solution}{exo:g9:ions:7}
Sulfate: 2 extra \omterm{def:g9:inside-the-atom:nucleus}{electrons}. Nitrate: 1.
\end{solution}

\begin{solution}{exo:g9:ions:8}
White with silver nitrate: \ce{Cl-}. Blue with sodium hydroxide:
\ce{Cu^{2+}}. Copper(II) chloride, \ce{CuCl2}.
\end{solution}

\begin{solution}{exo:g9:ions:9}
Four: one on each side (left, right, above, below).
\end{solution}

\begin{solution}{exo:g9:ions:10}
\ce{Fe^{2+}(aq) + 2OH-(aq) -> Fe(OH)2(s)}.
\end{solution}

\begin{solution}{exo:g9:ions:11}
Zinc chloride: sodium \omterm{def:g9:ions:ion}{ions} give no \omterm{def:g9:ions:precipitate}{precipitate} with sodium hydroxide,
zinc \omterm{def:g9:ions:ion}{ions} give a white one:
\ce{Zn^{2+}(aq) + 2OH-(aq) -> Zn(OH)2(s)}.
\end{solution}

\begin{solution}{exo:g9:ions:12}
\ce{CaCl2}: 2000 chloride \omterm{def:g9:ions:ion}{ions}. Charges: $1000 \times (+2) + 2000 \times
(-1) = 0$.
\end{solution}

\begin{solution}{pb:g9:ions:1}
\textbf{1.} Iron(II), \ce{Fe^{2+}}: \ce{Fe(OH)2}.

\textbf{2.} Iron(III), \ce{Fe^{3+}}: \ce{Fe(OH)3}.

\textbf{3.} \ce{Fe^{2+}} has one more: $26 - 2 = 24$ \omterm{def:g9:inside-the-atom:nucleus}{electrons}, against
$26 - 3 = 23$ for \ce{Fe^{3+}}.

\textbf{4.} \ce{Fe^{3+}(aq) + 3OH-(aq) -> Fe(OH)3(s)}.

\textbf{5.} It comes out holding iron(II) \omterm{def:g9:ions:ion}{ions}; in contact with the \omterm{def:g4:air-a-mixture-of-gases:air}{air}
they become iron(III) \omterm{def:g9:ions:ion}{ions}, which form the orange solid.

\textbf{6.} \ce{Fe(OH)3}: $(+3) + 3 \times (-1) = 0$.

\textbf{7.} Not as it comes out: the iron \omterm{def:g9:ions:ion}{ions} are dissolved and pass
through a filter. Once the orange solid has formed, yes: it is a
\omterm{def:g9:ions:precipitate}{precipitate}, held back by a filter.

\textbf{8.} No: a \omterm{def:g6:pure-substances-and-mixtures:mixture}{mixture} (water and dissolved \omterm{def:g9:ions:ion}{ions}). Freshly drawn and
clear, it is homogeneous.

\textbf{9.} $56 \times 1.67 \times 10^{-27} = \qty{9.35e-26}{kg}$.

\textbf{10.} $\qty{0.30}{mg} = \qty{0.30e-6}{kg} = \qty{3.0e-7}{kg}$.

\textbf{11.} An \omterm{def:g9:ions:ion}{ion} differs from the \omterm{def:g7:atoms-and-molecules:atom}{atom} by two or three \omterm{def:g9:inside-the-atom:nucleus}{electrons},
whose mass is negligible beside the \omterm{def:g9:inside-the-atom:nucleus}{nucleus}.

\textbf{12.} $3.0 \times 10^{-7} / 9.35 \times 10^{-26} \approx 3.2
\times 10^{18}$ iron \omterm{def:g9:ions:ion}{ions} per litre.
\end{solution}
